\documentclass[a0,portrait]{a0poster}

\usepackage{multicol} % This is so we can have multiple columns of text side-by-side
\columnsep=100pt % This is the amount of white space between the columns in the poster
\columnseprule=3pt % This is the thickness of the black line between the columns in the poster

\usepackage[svgnames]{xcolor} % Specify colors by their 'svgnames', for a full list of all colors available see here: http://www.latextemplates.com/svgnames-colors

\usepackage{times} % Use the times font
\usepackage{palatino} % Uncomment to use the Palatino font

\usepackage{graphicx} % Required for including images
\graphicspath{{figures/}} % Location of the graphics files
\usepackage{booktabs} % Top and bottom rules for table
\usepackage[font=small,labelfont=bf]{caption} % Required for specifying captions to tables and figures
\usepackage{amsfonts, amsmath, amsthm, amssymb} % For math fonts, symbols and environments
\usepackage{wrapfig} % Allows wrapping text around tables and figures
\usepackage{lipsum,adjustbox}
\usepackage[absolute,overlay]{textpos}
\usepackage{multirow}
\usepackage{url}
\usepackage{tikz}
\usepackage{tikz-dependency}
\usetikzlibrary{arrows.meta}
\newcommand{\parser}[1]{TUPA\textsubscript{#1}}
\captionsetup{labelformat=empty}

\begin{document}

% The header is divided into two boxes:
% The first is 75% wide and houses the title, subtitle, names, university/organization and contact information
% The second is 25% wide and houses a logo for your university/organization or a photo of you
% The widths of these boxes can be easily edited to accommodate your content as you see fit

\begin{center}
	\veryHuge \color{NavyBlue} \textbf{Broad-Coverage Transition-Based UCCA Parsing}
\end{center}
\vspace{-1cm}
\begin{minipage}[b]{.07\linewidth}
\includegraphics[width=\linewidth]{huji_logo.jpg}
\vspace{5mm}
\end{minipage}
\begin{minipage}[b]{.16\linewidth}
\includegraphics[width=\linewidth]{huji_banner.png}
\includegraphics[width=\linewidth]{cse_banner.jpg}
\vspace{1cm}
\end{minipage}
\hspace{1cm}
\begin{minipage}[b]{0.57\linewidth}
\LARGE \textbf{Daniel Hershcovich$^{1,2}$ \And Omri Abend$^2$ \And Ari Rappoport$^2$} \\[0.5cm] % Author(s)
\Large $^1$The Edmond and Lily Safra Center for Brain Sciences \\
  $^2$School of Computer Science and Engineering
  \setlength{\columnseprule}{0pt}
  \setlength\multicolsep{-20pt}
  \begin{multicols}{2}
  The Hebrew University of Jerusalem
  \large \texttt{\{danielh,oabend,arir\}@cs.huji.ac.il}
  \end{multicols}
\end{minipage}
\hfill
\begin{minipage}[b]{.13\linewidth}
\includegraphics[width=\linewidth]{elsc_logo.png}
\end{minipage}

\vspace{1cm} % A bit of extra whitespace between the header and poster content

%----------------------------------------------------------------------------------------


\color{Navy} % Navy color for the abstract

\begin{abstract}

  We present the first parser for UCCA, a
  cross-linguistically applicable framework for semantic
  representation, that builds on extensive
  typological work, and supports rapid annotation.
  UCCA poses a challenge for existing parsing techniques,
  as it exhibits reentrancy (resulting in DAG structures),
  discontinuous structures and non-terminal nodes corresponding
  to complex semantic units. To our knowledge, the conjunction
  of these formal properties is not supported by any existing parser.
  Our transition-based parser, using novel transition set
  and features, has value not just for UCCA parsing:
  its ability to handle more general graph structures will inform
  the development of parsers for other semantic DAG structures, 
  and in languages which frequently use discontinuous structures.

\end{abstract}


\begin{multicols}{2} % This is how many columns your poster will be broken into, a portrait poster is generally split into 2 columns


\color{Black} % SaddleBrown color for the introduction

\section*{Introduction}


\begin{wraptable}{R}{55mm}
  \vspace{-25mm}
  \begin{adjustbox}{margin=2mm,frame}
  \scalebox{.9}{
  \begin{tabular}{cl}
	  P & process \\
	  S & state \\
	  A & participant \\
	  L & linker \\
	  H & linked scene \\
	  C & center \\
	  E & elaborator \\
	  D & adverbial \\
	  R & relator \\
	  N & connector \\
	  U & punctuation \\
	  F & function unit \\
	  G & ground
  \end{tabular}
  }
  \end{adjustbox}
  \captionof{table}{Edge labels.}
\end{wraptable}

Universal Conceptual Cognitive Annotation \cite{abend2013universal},
is a cross-linguistically applicable semantic representation scheme.
UCCA has demonstrated applicability to multiple languages, including
English, French, German and Czech, support for rapid annotation,
and semantic stability in translation \cite{sulem2015conceptual}.
The scheme has proven useful for machine translation evaluation \cite{birch2016hume},
but its applicability has been so far limited by the absence of a UCCA parser.

Formally, a UCCA structure is a DAG whose leaves correspond to the tokens of
the text. Nodes (or {\it units}) either correspond to a terminal or
to several sub-units (not necessarily contiguous) jointly viewed as a
single entity according to some semantic or cognitive consideration.
Edges bear a category, indicating the role of the sub-unit in the relation
that the parent represents.

\begin{center}
  \begin{tikzpicture}[level distance=4cm, sibling distance=4cm, -{Latex[length=5mm]},
      every circle node/.append style={fill=black}]
    \node (ROOT) [circle] {}
      child {node (After) {After} edge from parent node[left] {L}}
      child {node (graduation) [circle] {}
      {
        child {node {graduation} edge from parent node[left] {P}}
      } edge from parent node[left] {H} }
      child {node {,} edge from parent node[right] {U}}
      child {node (moved) [circle] {}
      {
        child {node (John) {John} edge from parent node[left] {A}}
        child {node {moved} edge from parent node[left] {P}}
        child {node [circle] {}
        {
          child {node {to} edge from parent node[left] {R}}
          child {node {Paris} edge from parent node[right] {C}}
        } edge from parent node[right] {A} }
      } edge from parent node[right] {H} }
      ;
    \draw[dashed,-{Latex[length=5mm]}] (graduation) to node [auto] {A} (John);
  \end{tikzpicture}
  \captionof{figure}{Remote edge (dashed), resulting in ``John'' having two parents.}
  \begin{tikzpicture}[level distance=5cm, sibling distance=4cm, -{Latex[length=5mm]},
      every node/.append style={midway},
      every circle node/.append style={fill=black}]
    \node (ROOT) [circle] {}
      child {node {John} edge from parent node[left] {A}}
      child {node [circle] {}
      {
      	child {node {gave} edge from parent node[left] {C}}
      	child {node (everything) {everything} edge from parent[white]}
      	child {node {up} edge from parent node[right] {C}}
      } edge from parent node[right] {P} }
      ;
    \draw[bend right,-{Latex[length=5mm]}] (ROOT) to[out=-20, in=180] node [left] {A} (everything);
  \end{tikzpicture}
  \hspace{4cm}
  \begin{tikzpicture}[level distance=5cm, sibling distance=4cm, -{Latex[length=5mm]},
      every node/.append style={midway},
      every circle node/.append style={fill=black}]
    \node (ROOT) [circle] {}
      child {node [circle] {}
      {
        child {node {John} edge from parent node[left] {C}}
        child {node {and} edge from parent node[left] {N}}
        child {node {Mary} edge from parent node[right] {C}}
      } edge from parent node[left] {A} }
      child {node {went} edge from parent node[left] {P}}
      child {node {home} edge from parent node[right] {A}}
      ;
  \end{tikzpicture}
  \captionof{figure}{Discontinuous unit (``gave ... up''). \hspace{4cm}
  Coordination construction (``John and Mary'').}
  \captionof{figure}{UCCA structures demonstrating three structural properties.
	Pre-terminal nodes omitted for brevity.}
\end{center}



\section*{Transition-based UCCA Parsing}

\parser{} is a transition-based parser supporting the structural properties of UCCA.\footnote{The
corpora are available at \url{http://www.cs.huji.ac.il/~oabend/ucca.html},
and the parser code at \url{https://github.com/danielhers/ucca}.}
Transition-based parsers \cite{Nivre03anefficient} scan the text from start to end,
and create the parse incrementally by applying a \textit{transition}
at each step to the parser state,
defined using three data structures: a buffer $B$ of tokens and nodes to be processed,
a stack $S$ of nodes currently being processed,
and a graph $G=(V,E,\ell)$ of constructed nodes and edges,
where $V$ is the set of \emph{nodes}, $E$ is the set of \emph{edges},
and $\ell : E \to L$ is the \emph{label} function, $L$ being the set of possible labels.
Some of the states are marked as \textit{terminal}, meaning that $G$ is the final output.
A classifier is used at each step to select the next transition based on features
encoding the parser's current state.
During training, an oracle creates training instances for the classifier,
based on gold-standard annotations.

\begin{flushleft}
	\begin{tikzpicture}[level distance=27mm, sibling distance=4cm]
	\draw[xstep=5cm,ystep=1cm,color=gray] (-5.01,0) grid (5,1);
	\node[anchor=west] at (0,2) {$S$};
	\node[anchor=west] at (-5,.35) {After};
	\node[anchor=west] at (0,.35) {graduation};
	\draw[xstep=3cm,ystep=1cm,color=gray] (9,0) grid (21,1);
	\node[anchor=west] at (14,2) {$B$};
	\node[anchor=west] at (9,.35) {John};
	\node[anchor=west] at (12,.35) {moved};
	\node[anchor=west] at (15,.35) {to};
	\node[anchor=west] at (18,.35) {Paris};
	\node[anchor=west] at (25,3) {$G$};
	\node (ROOT) [fill=black, circle] at (26,1) {}
	  child {node  {After} edge from parent [-{Latex[length=5mm]}] node[left] {L}}
	  child {node [fill=black, circle] {}
	  {
	    child {node {graduation} edge from parent [-{Latex[length=5mm]}] node[right] {P}}
	  } edge from parent [-{Latex[length=5mm]}] node[right] {H} };
	\end{tikzpicture}
\end{flushleft}

\begin{center}
	\begin{adjustbox}{margin=3pt,frame}
	\begin{tabular}{llll|l|llll}
	\multicolumn{4}{c|}{\textbf{\small Before Transition}} & \textbf{\small Transition} & \multicolumn{4}{c}{\textbf{\small After Transition}} \\
	\textbf{\footnotesize Stack} & \textbf{\footnotesize Buffer} & \textbf{\footnotesize Nodes} & \textbf{\footnotesize Edges} & & \textbf{\footnotesize Stack} & \textbf{\footnotesize Buffer} & \textbf{\footnotesize Nodes} & \textbf{\footnotesize Edges} \\
	$S$ & $x \;|\; B$ & $V$ & $E$ & \textsc{Shift} & $S \;|\; x$ & $B$ & $V$ & $E$ \\
	$S \;|\; x$ & $B$ & $V$ & $E$ & \textsc{Reduce} & $S$ & $B$ & $V$ & $E$ \\
	$S \;|\; x$ & $B$ & $V$ & $E$ & \textsc{Node$_X$} & $S \;|\; x$ & $y \;|\; B$ & $V \cup \{ y \}$ & $E \cup \{ (y,x)_X \}$ \\
	$S \;|\; y,x$ & $B$ & $V$ & $E$ & \textsc{Left-Edge$_X$} & $S \;|\; y,x$ & $B$ & $V$ & $E \cup \{ (x,y)_X \}$ \\
	$S \;|\; x,y$ & $B$ & $V$ & $E$ & \textsc{Right-Edge$_X$} & $S \;|\; x,y$ & $B$ & $V$ & $E \cup \{ (x,y)_X \}$ \\
	$S \;|\; y,x$ & $B$ & $V$ & $E$ & \textsc{Left-Remote$_X$} & $S \;|\; y,x$ & $B$ & $V$ & $E \cup \{ (x,y)_X^* \}$ \\
	$S \;|\; x,y$ & $B$ & $V$ & $E$ & \textsc{Right-Remote$_X$} & $S \;|\; x,y$ & $B$ & $V$ & $E \cup \{ (x,y)_X^* \}$ \\
	$S \;|\; x,y$ & $B$ & $V$ & $E$ & \textsc{Swap} & $S \;|\; y$ & $x \;|\; B$ & $V$ & $E$ \\
	$[\mathrm{root}]$ & $\emptyset$ & $V$ & $E$ & \textsc{Finish} & $\emptyset$ & $\emptyset$ & $V$ & $E$ \\
	\end{tabular}
	\end{adjustbox}
	\captionof{table}{\parser{} transitions.
	  We write the stack with its top to the right and the buffer with its head to the left.\\
	  $(\cdot,\cdot)_X$ denotes a primary $X$-labeled edge, and $(\cdot,\cdot)_X^*$ a remote $X$-labeled edge.
	}
\end{center}


\paragraph{Classifier.}
We experiment with three classifiers:
(1) \parser{sparse}: perceptron with sparse features,
(2) \parser{dense}: perceptron with dense embedding features,
(3) \parser{NN}: feedforward neural network with embedding features.

For all classifiers, inference is performed greedily, i.e., without beam search.

\paragraph{Features.}
For the sparse perceptron-based parser (\parser{sparse}),
we use binary indicator features representing
the words, POS tags and existing edge labels related to the top four stack elements and the next
three buffer elements, in addition to their children and grandchildren in the graph.
We also use bi- and trigram features based on these values \cite{zhang2009transition,zhu2013fast},
features related to discontinuous nodes \cite{maier2015discontinuous},
and features representing existing edges and the number of parents and children a node has \cite{tokgoz2015transition}.
In addition, we use novel, UCCA-specific features:
the number of remote children a node has.

For the dense perceptron-based (\parser{dense}) and the NN-based parser (\parser{NN}),
we replace all binary features by a
concatenation of the vector embeddings of all the represented elements:
words, POS tags, edge labels and parser actions.

Finally, for all classifiers we employ a real-valued feature,
\textbf{ratio}, corresponding to the ratio between the number of terminals to number of nodes
in the graph $G$.
This novel feature serves as a regularizer for the creation of new nodes,
and should be beneficial for other transition-based constituency parsers too.




\section*{Experimental Setup}


\begin{center}
	\begin{tabular}{l|ccc|c}
	& \multicolumn{3}{c|}{Wiki} & 20K \\
	& \small Train & \small Dev & \small Test & Leagues \\
	\hline
	\# passages & 300 & 34 & 34 & 154 \\
	\# sentences & 4267 & 453 & 518 & 506 \\
	\hline
	\# nodes & 298,665 & 33,263 & 37,262 & 29,315 \\
	\% terminal & 42.95 & 43.62 & 42.89 & 42.09 \\
	\% non-term. & 58.30 & 57.46 & 58.31 & 60.01 \\
	\% discont. & 0.53 & 0.51 & 0.47 & 0.81 \\
	\% reentrant & 2.31 & 1.76 & 2.18 & 2.03 \\
	\hline
	\# edges & 287,381 & 32,015 & 35,846 & 27,749 \\
	\% primary & 98.29 & 98.81 & 98.75 & 97.73 \\
	\% remote & 1.71 & 1.19 & 1.25 & 2.27 \\
	\hline
	\multicolumn{3}{l}{\footnotesize Average per non-terminal node} \\
	\# children & 1.67 & 1.68 & 1.66 & 1.61 
	\end{tabular}
	\captionof{table}{Statistics of the \textit{Wiki} and \textit{20K Leagues} UCCA corpora.\\
	The latter was used for out-of-domain experiments.}
\end{center}


\paragraph{Evaluation.}
We define two simple measures for comparing UCCA structures.
Assume $G_p=(V_p,E_p,\ell_p)$ and $G_g=(V_g,E_g,\ell_g)$
are the predicted and gold-standard graphs over the same
sequence of terminals $W = \{w_1,\ldots,w_n\}$, respectively.
For an edge $e=(u,v)$ in either graph,
where $u$ is the parent and $v$ is the child, define its yield $y(e) \subseteq W$ as the
set of terminals in $W$ that are descendants of $v$.
Define the set of \textit{mutual edges} between $G_p$ and $G_g$:
\[
    M(G_p,G_g) =
    \left\{(e_1,e_2) \in E_p \times E_g \;|\;
    y(e_1) = y(e_2) \wedge \ell_p(e_1)=\ell_g(e_2)\right\}
\]
Labeled precision and recall are defined by dividing $|M(G_p,G_g)|$ by $|E_p|$ and $|E_g|$, respectively.
We report two variants of this measure: one where we consider only primary edges,
and another for remote edges.

\paragraph{Comparison to bilexical DAG parsers.}
As \parser{} is the first parser to support reentrancy, discontinuous units and
non-terminal nodes, no direct comparison with an existing parser is possible.
We instead compare \parser{} to bilexical DAG parsers,
which we view as the closest existing match (they support two of the properties,
but not non-terminal nodes).
Our approach is similar in spirit
to the \textit{tree approximation} approach used for dependency DAG parsing
\cite{agic2015semantic,fernandez2015parsing},
where dependency DAGs were converted into dependency trees
and then parsed by existing dependency tree parsers.

\begin{center}
	\begin{dependency}[theme = simple]
	\begin{deptext}[column sep=.7em,ampersand replacement=\^]
	After \^ graduation \^ , \^ John \^ moved \^ to \^ Paris \\
	\end{deptext}
	\depedge{2}{1}{L}
	\depedge{2}{3}{U}
	\depedge[dashed]{2}{4}{A}
	\depedge{5}{4}{A}
	\depedge{2}{5}{H}
	\depedge{7}{6}{R}
	\depedge{5}{7}{A}
	\end{dependency}\\
	\begin{dependency}[theme = simple]
	\begin{deptext}[column sep=.7em,ampersand replacement=\^]
	John \^ gave \^ everything \^ up \\
	\end{deptext}
	\depedge{1}{2}{A}
	\depedge{3}{2}{A}
	\depedge{4}{2}{C}
	\end{dependency}
	\hspace{1cm}
	\begin{dependency}[theme = simple]
	\begin{deptext}[column sep=.7em,ampersand replacement=\^]
	John \^ and \^ Mary \^ went \^ home \\
	\end{deptext}
	\depedge[edge start x offset=-6pt]{1}{4}{A}
	\depedge{2}{1}{N}
	\depedge{3}{1}{C}
	\depedge{5}{4}{A}
	\end{dependency}
	\captionof{figure}{Bilexical approximation.}
\end{center}

We first convert UCCA into bilexical dependencies, and train each parser on the resulting
training set.
We evaluate the trained parsers by applying them to the test set
and then reconstructing UCCA graphs, which are compared with the gold standard.
Upper bounds for the conversion-based method are computed by applying
the conversion and inverse conversion on the gold standard
graphs and comparing them to the original gold standard.

We evaluate two parsers in the bilexical approximation setting:
DAGParser \cite{ribeyre-villemontedelaclergerie-seddah:2014:SemEval}, the leading transition-based parser which participated in the SDP task of SemEval 2014 \cite{oepen2014semeval},
and the semantic variant of TurboParser \cite{almeida-martins:2015:SemEval},
a second-order graph-based parser that uses dual decomposition for optimization,
which obtained some of the best results in the SDP task of SemEval 2015 \cite{oepen2015semeval}. Default settings are used in all cases.





\section*{Results}

\parser{NN} obtains the highest scores in all metrics, surpassing the bilexical parsers
and the other classifiers:
	  
\begin{center}
	\begin{tabular}{l|ccc|ccc||ccc|ccc}
	& \multicolumn{6}{c||}{Wiki (in-domain)} & \multicolumn{6}{c}{20K Leagues (out-of-domain)} \\
	& \multicolumn{3}{c|}{Primary} & \multicolumn{3}{c||}{Remote}
	& \multicolumn{3}{c|}{Primary} & \multicolumn{3}{c}{Remote} \\
	& \textbf{LP} & \textbf{LR} & \textbf{LF} & \textbf{LP} & \textbf{LR} & \textbf{LF}
	& \textbf{LP} & \textbf{LR} & \textbf{LF} & \textbf{LP} & \textbf{LR} & \textbf{LF} \\
	\hline
	\multicolumn{4}{l}{\rule{0pt}{2ex} \footnotesize Bilexical Approximation} \\
	Upper Bound % on the test set and ood set only
	& 93.4 & 83.7 & 88.3 & 73.9 & 49.5 & 59.3
	& 93.5 & 83.5 & 88.2 & 66.7 & 31.6 & 42.9 \vspace{.1cm} \\
	DAGParser
	& 63.7 & 	56.1	 & 59.5	 & 0.8	 & 9.5	 &  1.4
	& 58	 & 49.8	 & 53.4 & -- & 0 & 0 \\
	TurboParser
	& 60.2	 & 47.4	 & 52.9	 & 2.2	 & 7.8	 &  3.4
	& 52.6	 & 39	 & 44.7	 & 100	 & 0.3	 & 0.6 \\
	\hline
	\multicolumn{4}{l}{\rule{0pt}{2ex} \footnotesize Direct Approach} \\
	\parser{sparse}
	& 64 & 55.6 & 59.5 & 16 & 11.6 & 13.4 
	& 60.6 & 53.9 & 57.1 & 20.2 & 10.3 & 13.6 \\
	\parser{dense} 
	& 55 & 54.8 & 54.9 & 15.2 & 16.9 & 16 
	& 54.8 & 55.2 & 55 & 6 & 3 & 4 \\
	\parser{NN}
	& {\bf 65} & {\bf 62.5} & {\bf 63.7} & {\bf 20.7} & {\bf 11.3} & {\bf 14.6}
	& {\bf 58.3} & {\bf 56.4} & {\bf 57.3} & {\bf 15.2} & {\bf 3.8} & {\bf 6}
	\end{tabular}
	\captionof{table}{Main experimental results in percents.}
\end{center}

Replacing the sparse representation with dense embeddings, \parser{dense} obtains
lower results than its sparse counterpart (although its recall on remote edges is
better in the in-domain setting).
Finally, using a feedforward NN, \parser{NN} manages to exploit the embeddings,
and obtains substantially higher scores in all measures, both
in the in-domain and out-of-domain settings.
Its performance in absolute terms, of 63.7\% F-score, is encouraging in light of
UCCA's inter-annotator agreement of 80--85\%
F-score on primary edges \cite{abend2013universal}.




\section*{Discussion}

\paragraph{Tree approximation.}
For completeness, we further explore a conversion-based approach that
lossily converts UCCA structures into trees. 
This approach resembles the bilexical approximation approach, but here it results in a simplified task for the underlying parser: parsing into tree structures.

Performance on remote edges is of pivotal importance in this investigation, which focuses on extending the class of graphs supported by statistical parsers. Nevertheless, we believe that exploring tree approximation methods, thus building on the maturity and cumulative experience with transition-based tree parsing, can inform the future development of DAG parsers in general and of UCCA parsers in particular.

We explore two conversion scenarios: one into (possibly discontinuous) constituency trees,
and one into bilexical dependency trees. In the first setting we experiment with \textsc{uparse}
\cite{maier-lichte:2016:DiscoNLP},
the only transition-based constituency parser, to our knowledge, able to parse trees with
discontinuous constituents.
In the second setting we use MaltParser \cite{nivre2007maltparser} with the
arc-eager transition set and the SVM classifier,
and the stack LSTM-based arc-standard parser \cite{dyer2015transition}.
Default settings are used in all cases.
\textsc{uparse} uses beam search by default,
with a beam size of 4, and the other parsers are greedy.
We also train the \parser{} variants without remote edge transitions, yielding a tree parser.

\begin{center}
\begin{tabular}{l|ccc}
& \textbf{LP} & \textbf{LR} & \textbf{LF} \\
\hline
\multicolumn{4}{l}{\rule{0pt}{2ex} \footnotesize Constituency Tree Approximation} \\
Upper Bound & 100 & 100 & 100 \vspace{.1cm} \\
\textsc{uparse} & 63 & 64.7 & 63.7 \\
\hline
\multicolumn{4}{l}{\rule{0pt}{2ex} \footnotesize Dependency Tree Approximation} \\
Upper Bound & 93.7 & 83.6 & 88.4 \vspace{.1cm} \\
MaltParser & 64.9 & 57.9 & 61 \\
LSTM Parser & {\bf 74.9} & {\bf 66.4} & {\bf 70.2} \\
\hline
\multicolumn{4}{l}{\rule{0pt}{2ex} \footnotesize Direct Tree Parsing} \\
\parser{sparse} $-$ \textsc{Remote} & 65.5 & 57.5 & 61.3 \\
\parser{dense} $-$ \textsc{Remote} & 57.2 & 57.3 & 57.2 \\
\parser{NN} $-$ \textsc{Remote} & 66.3 & 64.4 & 65.3 \\
\end{tabular}
\captionof{table}{
  Results of tree approximation experiment, in percents (on the \textit{Wiki} test set).
}
\end{center}

The conversion to constituency format only removes remote edges,
and thus obtains a perfect primary edge score.
The conversion to dependency format loses considerably more information, since
all non-terminal nodes have to be reconstructed by a simple rule-based inverse
conversion. Both conversions yield zero scores on remote edges,
since these are invariably removed when converting to trees.

The LSTM parser obtains the highest primary F-score,
with a considerable margin. Importantly, it obtains 9.2\% F-score
higher than the MaltParser,
despite being limited by the same approximation error upper bound,
and using a similar transition set.
This shows that the LSTM classifier has a significant contribution
to the accuracy of
the parser, suggesting that a recurrent neural network (RNN) model
may also improve \parser{}.
This is in line with the recent success of RNNs
(specifically, LSTMs and GRUs) in NLP tasks, including DAG-structured LSTM
\cite{zhu-sobhani-guo:2016:N16-1}.




\color{SaddleBrown} % SaddleBrown color for the conclusions to make them stand out

\section*{Conclusions}

We present \parser{}, the first parser for UCCA.
Evaluated in both in-domain and out-of-domain settings, we show that coupled with a
NN classifier, \parser{NN} surpasses bilexical DAG parsers on the task of UCCA parsing.
While much recent work focused on semantic parsing of different types,
the effectiveness of different parsing approaches for schemes with
different structural and semantic properties is not well-understood 
\cite{kuhlmann2016towards}.
We make a contribution to this literature by proposing the first grammarless parser that supports multiple parents, non-terminal nodes and discontinuous units, through detailed experiments with different architectures for the local classifier, and comparison both to bilexical approximations and to tree approximations.




\color{DarkSlateGray} % Set the color back to DarkSlateGray for the rest of the content

\section*{Future Work}

Future work will delve deeper in the exploration of
conversion-based parsing approaches, including different target
representations and more sophisticated conversion procedures \cite{kong-15},
in order to shed light on the commonalities and differences between
representations, suggesting ways for a data-driven design of semantic structures.
A parser for UCCA will enable using the framework for new tasks,
in addition to existing applications for evaluation of
machine translation \cite{birch2016hume}.
We believe that UCCA's merits in providing a cross-linguistically applicable,
broad-coverage annotation will support ongoing efforts to incorporate deeper
semantic structures into a variety of applications,
such as machine translation \cite{jones2012semantics}
and summarization \cite{liu2015toward}.




%\nocite{*} % Print all references regardless of whether they were cited in the poster or not
\bibliographystyle{plain} % Plain referencing style
\bibliography{references}


%----------------------------------------------------------------------------------------

\end{multicols}
\end{document}
